\begin{frame}[label=q5b2]{5b.2 · DataFrame: ordenamiento decreciente}
  \small
  \pregunta{¿Cuáles son los 15 distritos con más personas reportadas durante 2019--2024?}
  \entrada{93\,831 filas. \texttt{departamento}, \texttt{provincia}, \texttt{distrito}, \texttt{cantidad}.}
  \flujo{Todos los registros}{\texttt{groupBy}\\Departamento, provincia, distrito}{\texttt{SUM(cantidad)}\\\texttt{orderBy(desc)}}{\texttt{limit(15)}\\15 distritos}
  \begin{columns}[T,onlytextwidth]
    \column{0.55\textwidth}
    \textbf{Código Scala original (fragmento)}
    \codigo{sql/05b_2_top_distritos_desc}{10}{16}
    \column{0.42\textwidth}
    \textbf{Salida}\\[5pt]
    {\footnotesize
\begin{tabular}{@{}ll@{}}
\toprule
Distrito & Personas\\
\midrule
San Juan de Lurigancho & 4\,775\\
Lima & 3\,153\\
San Martín de Porres & 2\,646\\
Callao & 2\,642\\
\bottomrule
\end{tabular}
    }
  \end{columns}
  \respuesta{San Juan de Lurigancho encabeza el acumulado con 4\,775. Las tres claves identifican la ubicación; el orden usa personas, no filas.}
  \origen{sql}{05b_2_top_distritos_desc}
  \note{20 s. Leer el primer lugar y mostrar orderBy(desc) seguido de limit.}
\end{frame}
